\documentclass[11pt]{article}
\usepackage[margin=1in]{geometry}
\usepackage{amsmath,amssymb,amsthm,booktabs,longtable,graphicx}
\usepackage[colorlinks=true,urlcolor=blue,linkcolor=blue]{hyperref}
\newtheorem{proposition}{Proposition}
\newtheorem{corollary}{Corollary}
\newcommand{\M}{\mathcal M}
\newcommand{\PrA}{\Pr_A}
\newcommand{\IdentifierEvaluations}{960}

\title{Exact randomization inference with uncertain record identities}
\author{Research draft}
\date{October 2026}
\begin{document}
\maketitle
\begin{abstract}
An experimental conclusion can depend on how outcome records are assigned to experimental
units. We develop a sensitivity analysis that retains candidate links, one-to-one ownership,
and a bound on changed accepted identities. Under complete randomization and full overlap,
every feasible outcome permutation has the same sharp-null reference distribution.
The largest Fisher p-value therefore follows from the smallest attainable absolute
contrast; the smallest p-value follows from the linear contrast endpoints. Integer
optimization returns the linkages attaining both extremes. We also report the minimum
identity-change budget that removes rejection and give a closed-form binary comparator
that uses accuracy information without candidate restrictions. Exhaustive checks establish
agreement with direct enumeration, and an identifier-based simulation measures what the
candidate graph adds. In these synthetic settings, gains are modest and substantial
power loss remains. The method specializes established matching-set and sensitivity
principles. Its inferential guarantee depends on the true linkage belonging to the stated
set; it does not establish that membership from name scores alone.
\end{abstract}

\section{The question and the contribution}
A researcher has randomized treatment assignments in one file and outcomes in another.
Several outcome records may plausibly belong to each experimental unit. The question is
whether the sharp no-effect null can be rejected under every permitted assignment of
records to units. A linkage accuracy rate restricts how many accepted identities may be
wrong. Candidate links additionally restrict which outcomes can move between which units.
These restrictions can change the answer even when the accepted matching and error budget
are held fixed.

Our contribution is a specified randomization sensitivity procedure over record identities.
It computes the exact p-value endpoints for a stated graph and error budget, returns the
attaining identities, and identifies the budget at which rejection first fails. Its useful
specialization is the combination of the experimental assignment law with conservation of
whole donor records. A corrected outcome vector that cannot arise from a feasible matching
is excluded. The algorithm does not require probabilities for individual links.

Optimization of statistics over matching confidence sets is established
\cite{hall}; integer-programming sensitivity to binary outcome misclassification is also
established \cite{heng}. Permutation-set inference under a Gaussian regression model uses
a different source of randomness \cite{ota}. The result here supplies the full-bijection
Fisher reference law and its reduction to attainable contrasts. It uses standard assignment
constraints, absolute-value optimization, and confidence-set maximization. The contribution
is this causal specialization and its implementation, rather than a new optimization or
coverage principle. Linkage-induced causal bias and loss of power already motivate a
substantial literature \cite{wortman,tahamont}.

\section{Experimental units and feasible identities}
Fix $N$ units with potential outcomes $Y_i(0),Y_i(1)$. Exactly $m$ receive treatment,
uniformly over the $\binom Nm$ assignments; write $c=N-m$. Assignment $Z$ is observed without
error. The donor file contains exactly one realized outcome for every unit and no other
outcomes. Its rows $y_1,\ldots,y_N$ have unknown true ownership $\pi^*$, a bijection.
Thus $Y_i=Y_i(Z_i)=y_{\pi^*(i)}$.

Let $G$ be a bipartite candidate graph and $\pi_0$ an accepted bijection. For integer $K$,
define
\[
 \M(G,K)=\left\{\pi:\pi\text{ bijective},\ (i,\pi(i))\in G\ \forall i,
 \quad \sum_i1\{\pi(i)\ne\pi_0(i)\}\le K\right\}.
\]
The analysis requires a nonempty set. Unless stated otherwise, it assumes $\pi^*\in\M(G,K)$.
A donor supplies its complete record: treating each field as independently reassignable
would describe a different uncertainty set. The present test concerns known assignment
and uncertain outcomes; it is not a general correction for jointly mislinked treatment,
covariates, and outcomes.

For any feasible $\pi$, the observed contrast is
\[
 T(\pi)=\sum_iw_i y_{\pi(i)},\qquad
 w_i=\frac{Z_i}{m}-\frac{1-Z_i}{c}.
\]
Let $L=\min_{\pi\in\M}T(\pi)$ and $U=\max_{\pi\in\M}T(\pi)$.
These are extrema of reconstructed estimates. They are not confidence limits for the
finite-population average effect $N^{-1}\sum_i\{Y_i(1)-Y_i(0)\}$.

\section{Exact p-value endpoints}
We test Fisher's sharp null $Y_i(1)=Y_i(0)$ for every unit. For a uniform size-$m$
assignment $A$, define the common tail function
\[
 H(s)=\PrA\left(\left|\sum_i\left\{\frac{A_i}{m}-\frac{1-A_i}{c}\right\}y_i\right|
 \ge s\right).
\]
\begin{proposition}[Common reference law]
Under complete randomization, the distribution of the absolute contrast is identical
for every full permutation of the donor outcomes. Consequently, if
$d=\min_{\pi\in\M}|T(\pi)|$ and $D=\max(|L|,|U|)$, then
\[
 p_{\max}=\max_{\pi\in\M}p(\pi)=H(d),\qquad
 p_{\min}=\min_{\pi\in\M}p(\pi)=H(D).
\]
Both endpoints are attained by feasible linkages.
\end{proposition}
\begin{proof}
A permutation of unit labels maps the collection of size-$m$ treatment subsets bijectively
onto itself. Each subset has the same probability, so permuting the outcome vector
preserves its randomization law. The function $H$ is nonincreasing, and the nonempty
finite feasible set attains the minimum and maximum absolute contrasts. The maximum
absolute contrast is attained at a linear endpoint.
\end{proof}

The maximum p-value answers whether every permitted linkage rejects. The minimum
p-value shows how favorable a permitted reconstruction could be. Selecting the linkage
with the smallest p-value after observing outcomes does not yield a valid robust test.
The reference law uses the full experimental assignment support; it does not condition
assignment on the observed candidate graph.

For binary outcomes with $s_y$ successes, treated successes follow
$Q\sim\operatorname{Hypergeometric}(N,s_y,m)$, and
\[
 H(s)=\sum_{q=\max(0,m-N+s_y)}^{\min(m,s_y)}
 \frac{\binom{s_y}{q}\binom{N-s_y}{m-q}}{\binom Nm}
 1\left\{\left|\frac q m-\frac{s_y-q}{c}\right|\ge s\right\}.
\]
Any two-valued outcome has the same calculation after affine recoding. For general
outcomes, the implementation enumerates the assignment distribution subject to an
explicit size limit. It optimizes over linkages without enumerating them. This does
not imply that integer optimization is inexpensive for arbitrary graphs or outcomes.

\subsection{Inference when the graph can miss truth}
\begin{proposition}[Containment-adjusted validity]
Suppose $\Pr\{\pi^*\notin\M(G,K)\}\le\delta$ under every null population under
consideration. Then $p_{\mathrm{rob}}=\min(1,p_{\max}+\delta)$ satisfies
$\Pr(p_{\mathrm{rob}}\le\alpha)\le\alpha$.
\end{proposition}
\begin{proof}
On containment, $p_{\max}\ge p(\pi^*)$ pointwise. For $\delta\le\alpha<1$, rejection
therefore implies either containment failure or $p(\pi^*)\le\alpha-\delta$.
The first event has probability at most $\delta$ and the second at most
$\alpha-\delta$ by oracle randomization validity. If $\alpha<\delta$, rejection
is impossible; $\alpha=1$ is immediate.
\end{proof}
This is the established confidence-set maximization argument \cite{berger}.
The graph may be data dependent; independence between graph construction and outcomes
is not required by the proof. What is required is the joint containment bound under
the relevant repeated experiment. Pairwise candidate recall is not the probability
that the entire true matching survives. The software accepts $\delta$; it does not
estimate or certify it.

\section{Optimization and an identity breakdown budget}
Use binary variables $x_{ij}$ for allowed edges. Require each row and column to sum to
one and $N-\sum_i x_{i,\pi_0(i)}\le K$. Minimizing and maximizing
$\sum_{ij}w_i y_j x_{ij}$ gives $L,U$. For $d$, add $t\ge0$ and minimize $t$ subject to
\[
 -t\le\sum_{ij}w_i y_jx_{ij}\le t.
\]
For two-valued outcomes, we optimize the equivalent integer objective
$|Nq-ms_y|$ and bound it below by its minimum over hypergeometric support. This
prevents a search for zero when the outcome lattice excludes it. The returned
assignments are checked for feasibility and agreement with the objectives.
An interval $[L,U]$ spanning zero does not guarantee an attainable zero. For example,
the only contrasts might be $-2$ and $2$. Substituting the distance of $[L,U]$ from
zero would discard the discreteness that the integer program preserves.

An identity swap between anchors $i,k$ changes the contrast by
$(w_i-w_k)(y_{\pi(k)}-y_{\pi(i)})$. Within-arm swaps preserve the unadjusted contrast;
cross-arm swaps matter according to the outcome difference. Arbitrary matching
differences decompose into alternating cycles. This explains why error counts and
candidate counts alone do not determine uncertainty. These are standard assignment
properties applied to the causal contrast, consistent with earlier observations about
harmless within-subclass linkage errors \cite{wortman}.

For breakdown analysis require the accepted map itself to satisfy the graph
constraints, so the zero-budget set is nonempty. Define the identity breakdown budget by
\[
 K_{\mathrm{break}}=\min\{K:p_{\mathrm{rob}}(G,K)>\alpha\},
\]
where rejection uses $p\le\alpha$. The output is zero if the containment-adjusted accepted p-value already fails
to reject and infinity if every graph-feasible map rejects. Because feasible sets grow
with $K$, $p_{\mathrm{rob}}$ is nondecreasing, so binary search finds the first budget.
It returns a linkage demonstrating the loss of rejection. This is an identity-constrained
counterpart of warning accuracy for outcome misclassification \cite{heng}, not a separate
claim to invent breakdown analysis.

\subsection{A comparator that conserves donors}
The accuracy-only comparison retains the same $y,\pi_0,K$ and one-to-one constraint,
but permits every anchor--donor edge. It does not invent replacement outcomes.
\begin{proposition}[Binary accuracy-only support]
For binary outcomes, let $q_0$ be accepted treated successes and $s_y$ total successes.
Under the complete graph and at most $K$ changed identities, attainable treated-success
counts are exactly the integers in
\[
 \left[\max\{0,m-N+s_y,q_0-\lfloor K/2\rfloor\},\quad
 \min\{m,s_y,q_0+\lfloor K/2\rfloor\}\right].
\]
\end{proposition}
\begin{proof}
Changing treated successes by $h$ requires at least $|h|$ changes of treated outcomes
and $|h|$ opposite changes of control outcomes. Hence at least $2|h|$ identities must
change. Disjoint swaps of opposite-valued donors across arms attain this lower bound
for every count within the natural support.
\end{proof}
Evaluating $H$ at these counts gives an exact comparator without a complete-graph
integer program. Arbitrary candidate graphs can leave gaps in attainable counts;
the interval formula is not transferred to them. Equal-valued donors also cannot be
collapsed if their allowed neighbors or accepted-identity costs differ.

\section{Validation and the information gained from candidates}
The code is checked against exhaustive linkages and assignments in small problems,
including unequal arms, off-graph budgets, ties, disconnected components, and shared
identity budgets. Tests compare p-value endpoints and their witnesses, rather than
only solver objective values. Further checks establish scale and location invariance
and compare breakdown budgets with exhaustive search. Outcomes are centered and
scaled before optimization and tail comparisons, then reported contrasts are returned
in their original units. This prevents an absolute numerical tolerance from making
small-unit outcomes appear identically zero. Exactness refers to the finite-sample
probability and global optimization targets; solver and probability calculations
still use floating-point tolerances, not arbitrary-precision arithmetic.

\subsection{A controlled comparison}
A deliberately constructed example keeps the accepted map, its identity accuracy,
and the number of candidates per record fixed while changing which cross-arm
exchanges are possible. Table~\ref{tab:example} shows different robust conclusions.
These are controlled information sets, not estimated graphs.
\begin{table}[htbp]\centering
\caption{Same accepted identities and accuracy information, different candidate graphs.
The controlled example fixes eight units, four treated, and an at-most-two identity-error
budget. Both restricted graphs offer two candidates per unit. P-values are exact
two-sided Fisher probabilities.}\label{tab:example}
\begin{tabular}{lrrr}\toprule
Information & Lower contrast & Upper contrast & Maximum p \\
\midrule
Accuracy only & 0.50 & 1.00 & 0.486 \\
Within-arm candidates & 1.00 & 1.00 & 0.029 \\
Cross-arm candidates & 0.50 & 1.00 & 0.486 \\
\bottomrule\end{tabular}\end{table}

Enumerating every assignment in the accompanying design checks size and power without
Monte Carlo uncertainty. Table~\ref{tab:design} retains the cases where the graph gains
little or substantial power is still lost. A narrow estimate range by itself does not
remove experimental uncertainty.
\begin{table}[htbp]\centering
\caption{Exact rejection probabilities in the eight-unit design at the 5\% level.
All 70 assignments are enumerated, so there is no Monte Carlo error.
The zero-effect row measures sharp-null size; other rows measure power under constant
additive effects. Each method uses the same accepted linkage and error budget.}
\label{tab:design}
\begin{tabular}{rrrrr}\toprule
Effect & Oracle & Accuracy only & First graph & Second graph \\
\midrule
0 & 0.029 & 0.000 & 0.029 & 0.000 \\
1 & 0.014 & 0.000 & 0.014 & 0.000 \\
2 & 1.000 & 0.000 & 0.086 & 0.086 \\
4 & 1.000 & 0.000 & 0.086 & 0.086 \\
\bottomrule\end{tabular}\end{table}


\subsection{Candidates generated from identifiers}
The second comparison contains \IdentifierEvaluations{} evaluations and generates
identifiers before treatment and outcomes. Records
share a known geographic group and a short, potentially duplicated identifier. At most
one identifier character is corrupted; candidate links require the same group and at
most one character difference. True-link containment follows from this stated error
support, without adding missing true edges after graph construction. A minimum-distance
one-to-one matching, with random outcome-independent tie breaking, supplies the accepted map.
Donor rows are shuffled so that row position does not reveal identity. There are
ten geographic groups in each population; increasing sample size also increases
the records per group and candidate ambiguity. This is not a scaling experiment
that holds linkage difficulty fixed.

The graph and accuracy-only procedures receive the same accepted map and the same
\emph{oracle-known} number of incorrect accepted identities. This isolates the value
of candidate restrictions; it is not a method for learning an accuracy budget from
unvalidated records. Simulations cross two sample sizes, two corruption probabilities,
uniform versus geographically varying baseline risks, and a sharp null plus two
positive treatment-response increments. Complete randomization is independent of
identifiers. Positive increments need not equal the average treatment effect because
success probabilities are capped at one. Every scenario is retained, including weak
gains and failures. A solver time limit gives a conservative nonrejection on failed
optimizations; those runs are recorded as unresolved, not exact p-value endpoints. The protocol records the previously inspected controlled examples
and separates timing pilots from final simulation seeds.
\begin{center}\small
\setlength{\LTcapwidth}{\textwidth}
\begin{longtable}{rrlrrrrr}
\caption{Identifier-generated candidates: rejection frequencies and paired gains.
Each row uses 40 replications at the 5\% level.
Corruption changes at most one identifier
character; geographic baseline risks vary by the known identifier group. Each comparison
uses the oracle-known accepted identity-error count. Rejection rates are proportions;
the last column is graph minus accuracy-only rejection, with Monte Carlo standard error.
Unresolved optimizations conservatively do not reject; failures are reported in the text.
The effect column is a probability increment before capping.}\label{tab:identifier}\\
\toprule
$N$ & Noise & Baseline & Effect & Oracle & Accuracy & Graph & Gain (SE) \\
\midrule\endfirsthead
\toprule $N$ & Noise & Baseline & Effect & Oracle & Accuracy & Graph & Gain (SE) \\
\midrule\endhead
40 & 0.25 & Geographic & 0.00 & 0.025 & 0.000 & 0.000 & 0.000 (0.000) \\
40 & 0.25 & Geographic & 0.25 & 0.175 & 0.000 & 0.025 & 0.025 (0.025) \\
40 & 0.25 & Geographic & 0.50 & 0.575 & 0.075 & 0.100 & 0.025 (0.025) \\
40 & 0.75 & Geographic & 0.00 & 0.000 & 0.000 & 0.000 & 0.000 (0.000) \\
40 & 0.75 & Geographic & 0.25 & 0.100 & 0.000 & 0.000 & 0.000 (0.000) \\
40 & 0.75 & Geographic & 0.50 & 0.575 & 0.000 & 0.100 & 0.100 (0.048) \\
40 & 0.25 & Uniform & 0.00 & 0.025 & 0.000 & 0.000 & 0.000 (0.000) \\
40 & 0.25 & Uniform & 0.25 & 0.225 & 0.050 & 0.050 & 0.000 (0.000) \\
40 & 0.25 & Uniform & 0.50 & 0.800 & 0.200 & 0.200 & 0.000 (0.000) \\
40 & 0.75 & Uniform & 0.00 & 0.025 & 0.000 & 0.000 & 0.000 (0.000) \\
40 & 0.75 & Uniform & 0.25 & 0.125 & 0.000 & 0.000 & 0.000 (0.000) \\
40 & 0.75 & Uniform & 0.50 & 0.850 & 0.025 & 0.050 & 0.025 (0.025) \\
80 & 0.25 & Geographic & 0.00 & 0.025 & 0.000 & 0.000 & 0.000 (0.000) \\
80 & 0.25 & Geographic & 0.25 & 0.500 & 0.000 & 0.000 & 0.000 (0.000) \\
80 & 0.25 & Geographic & 0.50 & 0.925 & 0.000 & 0.000 & 0.000 (0.000) \\
80 & 0.75 & Geographic & 0.00 & 0.075 & 0.000 & 0.000 & 0.000 (0.000) \\
80 & 0.75 & Geographic & 0.25 & 0.625 & 0.000 & 0.000 & 0.000 (0.000) \\
80 & 0.75 & Geographic & 0.50 & 0.950 & 0.000 & 0.000 & 0.000 (0.000) \\
80 & 0.25 & Uniform & 0.00 & 0.075 & 0.000 & 0.000 & 0.000 (0.000) \\
80 & 0.25 & Uniform & 0.25 & 0.500 & 0.000 & 0.000 & 0.000 (0.000) \\
80 & 0.25 & Uniform & 0.50 & 1.000 & 0.000 & 0.000 & 0.000 (0.000) \\
80 & 0.75 & Uniform & 0.00 & 0.025 & 0.000 & 0.000 & 0.000 (0.000) \\
80 & 0.75 & Uniform & 0.25 & 0.450 & 0.000 & 0.000 & 0.000 (0.000) \\
80 & 0.75 & Uniform & 0.50 & 0.975 & 0.000 & 0.000 & 0.000 (0.000) \\
\bottomrule\end{longtable}\end{center}

At the larger probability increment, the graph procedure rejects in 18 of 320 evaluations, compared with 12 for accuracy alone and 266 with known identities. The gains are modest and substantial power loss remains.
Across the positive-increment scenarios, the graph's paired rejection gain ranges from 0.00 to 0.10. Mean candidate counts range from 2.3 to 4.1 across scenarios; mean accepted identity accuracy ranges from 0.41 to 0.87.
Scenario median graph-test runtimes range from 0.001 to 0.010 seconds; the largest scenario 95th-percentile runtime is 3.808 seconds. These timings describe the tested graphs and hardware, not worst-case complexity.
Of 960 evaluations, 14 did not return certified optima under the per-solver time limit and received the conservative nonrejection fallback.

The intervals and standard errors in this comparison concern Monte Carlo uncertainty
across synthetic populations and assignments. A zero estimated Monte Carlo standard
error means all sampled paired differences were zero; it does not establish a zero
population difference. Completed individual tests use exact
hypergeometric reference probabilities. Synthetic name errors with known support and
an oracle error budget do not establish performance for actual administrative linkage.

\section{Scope and empirical use}
The method applies when the experiment is completely randomized, assignment is known,
and both files cover exactly the same units once. Blocked assignment can give different
reference laws when a linkage moves outcomes across blocks. Partial overlap or extra
donors can change the outcome multiset. Neither situation inherits the common-law
shortcut. Covariate-dependent regression weights and linked treatment labels require
separate derivations. The present null is no effect for every unit; the test is not
an exact test of a zero average effect with heterogeneous individual effects.

The Rajasthan reservation--MNREGA exercise remains a separate sensitivity application.
It assumes the reservation lottery requested for that analysis, uses a stratified design
and a rectangular candidate graph, and evaluates feasible reconstructions against their
own simulated assignment laws. Those outputs are not global exact p-value endpoints
from the theorem here. Its audit shows that dramatic coefficient movements depend on
weak GP-name candidates admitted by concatenated geographic names. Tighter candidate
rules greatly reduce those movements, but do not independently validate the remaining
identities. It therefore motivates scrutiny of feasible sets rather than validating
the present method's practical gain.

For an application, supply an independently justified candidate set and error budget,
compute the robust p-value and attaining map, and inspect the identity changes that
matter. If whole-set coverage is uncertain, report sensitivity to omissions or a
supported containment adjustment. The immediate empirical requirement is a validation
sample or benchmark with independently known identities. A small maximized p-value
then has a clear interpretation: every identity reconstruction retained by the stated
restrictions rejects the sharp no-effect null under the experimental design.

\begingroup\small
\begin{thebibliography}{9}
\bibitem{hall} Hall, R., and S. Fienberg (2012).
Valid Statistical Inference on Automatically Matched Files.
\url{https://www.cs.cmu.edu/~rjhall/linkage.pdf}.
\bibitem{heng} Heng, S., and P. A. Shaw (2025).
Sensitivity Analysis for Binary Outcome Misclassification in Randomization Tests via
Integer Programming. \url{https://arxiv.org/abs/2201.03111}.
\bibitem{ota} Ota, H., and M. Imaizumi (2026).
Finite-Sample Inference for Sparsely Permuted Linear Regression.
\url{https://arxiv.org/abs/2601.14872}.
\bibitem{berger} Berger, R. L., and D. D. Boos (1994).
P Values Maximized Over a Confidence Set for the Nuisance Parameter.
Journal of the American Statistical Association 89:1012--1016.
\url{https://doi.org/10.1080/01621459.1994.10476836}.
\bibitem{wortman} Wortman, J. H., and J. P. Reiter (2018).
Simultaneous Record Linkage and Causal Inference with Propensity Score Subclassification.
\url{https://arxiv.org/abs/1709.03631}.
\bibitem{tahamont} Tahamont, S., et al. (2021).
Dude, Where's My Treatment Effect? Errors in Administrative Data Linking and the
Destruction of Statistical Power in Randomized Experiments.
\url{https://doi.org/10.1007/s10940-020-09461-x}.
\end{thebibliography}
\endgroup
\end{document}
